\scanpage{118}
\mspage{112}
% archive sheet number: 117

L'axiome des supports ne concerne directement que les données $\unlhd$, $|\unsure{0|o}|$ et $\unsure{\mathscr{F}|\mathcal{F}|F}\subset\unsure{\mathscr{P}|\mathcal{P}|P}(\unsure{\mathcal{A}|\mathscr{A}|\mathbb{A}})$ sur $A$, et ne fait pas intervenir encore $\overset{\circ}{\ll}$ ou $\ll$. Je \unsure{convient|conviens} de le renforcer ainsi.

Considérons dans un \unsure{qtopos|topos} $\Sigma$ toutes les \deleted{familles $(e_{\ast},I)=(e_i)_{i\in I}$}\inserted{parties} \deleted{$e_{\ast}=\{e_i\}_{i\in\unsure{J|I}}$} $A$, finies d'objets \deleted{\unsure{non}} nuls mutuellement disjoints. On peut les considérer comme des objets de $\operatorname{Figspat}(\Sigma)$, pour la relation d'ordre discret (disons) sur $\deleted{\unsure{\Sigma|e_{\ast}}}A$. Nous allons écrire, pour
\[
A,B\in\operatorname{Figspatdisc}(\Sigma),
\]
\[
\deleted{\unsure{A}}B\ll\deleted{\unsure{B}}A\quad\Longleftrightarrow\quad
\forall\deleted{\unsure{a}}b\in\deleted{\unsure{A}}B,\quad
\exists\deleted{\unsure{b}}a\in\deleted{\unsure{B}}A\text{ avec }b\leq a
\]
(\uline{NB} c'est aussi la relation \inserted{induite par} $\ll$ de $\operatorname{Figspat}$)

\deleted{Car} Pour un $\deleted{\unsure{b}}a\in\deleted{\unsure{B}}A$ donné, considérons l'ens $\deleted{\unsure{A_b}}B_a=\{b\in\deleted{\unsure{A}}B\mid b\leq a\}$, formé des éléments de $\deleted{\unsure{A}}B$ (mut. disj.) qui sont $\leq a$. Soit $\deleted{\unsure{b'}}a'=\operatorname{Sup}\deleted{\unsure{A_b}}B_a\leq\deleted{\unsure{b}}a$, et $\deleted{\unsure{b''}}a''=b\wedge\unsure{\complement|C}b'$.
